\section{Conclusion}\label{sec:conclusion}

We investigated whether lightweight residual gating can recover goal attention in multi-turn dialogue. ThreadGate applies per-layer sigmoid gates to the residual stream, trained jointly with an auxiliary Goal Attention Ratio (GAR) loss that targets attention to system prompt goal tokens. Training succeeds: teacher-forced GAR reaches 0.49 across all fine-tuned variants. The failure emerges at inference time. Under autoregressive generation, GAR drops to 0.016 for ThreadGate, a 30$\times$ gap that the gating mechanism cannot bridge. A four-way ablation (ThreadGate, Uniform Gate, LoRA-only, Base) confirms that gating contributes zero measurable benefit, with GAR-ours varying by only 0.75 percentage points across fine-tuned variants on held-out data. A PCGrad diagnostic exclusion experiment rules out gradient direction conflict between the NTP and GAR objectives as the cause: eliminating direction conflict does not rescue GAR convergence. The root cause is exposure bias. Attention patterns optimized under teacher-forced prefixes do not transfer to model-generated text, and the primary NTP objective's self-recovery from exposure bias does not extend to auxiliary attention losses. These results indicate that auxiliary attention losses face a fundamental generalization barrier under teacher-forcing. Future work should explore inference-time attention intervention or exposure bias mitigation strategies that operate on the autoregressive distribution directly.